\chapter{Introduction}
\label{ch:intro}

{\meq}, the Maryland Equilibrium Solver, computes axisymmetric plasma
equilibria. It takes a plasma boundary and the profiles that drive the plasma,
and returns the poloidal flux function whose level sets are the flux surfaces
--- together with the field built from its gradient, the surfaces themselves as
curves, and the flux-surface averages a transport code reads off them.

This manual is meant to be read through. It sets out what the code is doing and
why that is correct: the operator and the weights that distinguish it from a
Laplacian, the first-order system and the reason the flux is solved for rather
than differentiated, what hybridization buys and what it costs, the obligation
Newton places on every source term, the transfer technique that lets a straight
mesh carry a curved plasma boundary, and the machinery that turns a
two-dimensional flux function into one-dimensional geometry. The online
documentation is the companion reference: it is organised for searching, and it
is where the key-by-key description of a configuration file, the recipes and the
API listing live. Where a topic belongs in both, this manual carries the physics
and the argument, and points onward for the catalogue.

\section{The problem}

An axisymmetric magnetohydrodynamic equilibrium in cylindrical coordinates
$(r, \varphi, z)$ is described entirely by one scalar: the poloidal flux per
radian, $\psi(r,z)$. Force balance --- the statement that the Lorentz force
balances the pressure gradient --- reduces to a single semi-linear elliptic
equation for it, the Grad--Shafranov equation,
%
\begin{equation}
	-\gradbar \cdot \left( \frac{1}{r}\, \gradbar \psi \right)
		= \frac{F(r,z,\psi)}{r}
	\qquad \text{in } \Omega \subset \mathbb{R}^{2},
	\label{eq:gs}
\end{equation}
%
in which the source
%
\begin{equation}
	F(r,z,\psi) := \mu_{0}\, r^{2}\, \frac{\mathrm{d}p}{\mathrm{d}\psi}
		+ g\, \frac{\mathrm{d}g}{\mathrm{d}\psi}
	\label{eq:source}
\end{equation}
%
is built from the plasma pressure $p(\psi)$ and the toroidal field function
$g(\psi)/r$. Both are user input, and it is their dependence on $\psi$ that
makes the problem semi-linear rather than linear.

\subsection{Notation, and the weights}

Two pieces of notation carry the whole content of \eqref{eq:gs} and are worth
stating before anything else.

The symbol $\gradbar := (\partial_{r}, \partial_{z})$ acts formally like a
vector of partial derivatives. \emph{It is not the cylindrical gradient}, and
the operator in \eqref{eq:gs} is not a cylindrical divergence of a cylindrical
gradient: writing it that way gives a different equation. The distinction is
exactly what the $1/r$ inside and the $r$ implied by the right-hand side are
for. Expanded,
%
\begin{equation}
	\dstar \psi
		:= r \, \gradbar \cdot \left( \frac{1}{r} \gradbar \psi \right)
		 = \frac{\partial^{2}\psi}{\partial r^{2}}
		 - \frac{1}{r} \frac{\partial \psi}{\partial r}
		 + \frac{\partial^{2}\psi}{\partial z^{2}},
	\label{eq:deltastar}
\end{equation}
%
which is the Grad--Shafranov operator, and \eqref{eq:gs} reads
$\dstar \psi = -F$. Note the \emph{minus} on the first-derivative term: an
axisymmetric Laplacian carries a plus there. Every sign convention in
{\meq} descends from \eqref{eq:gs} as written.

\emph{Fixed boundary} means the plasma boundary $\Gamma$ is given rather than
found. Taking it to be the level set $\psi = 0$ makes \eqref{eq:gs} an interior
Dirichlet problem,
%
\begin{equation}
	\psi = 0 \qquad \text{on } \Gamma = \partial\Omega .
\end{equation}
%
That is the problem {\meq} solves.

\section{Three commitments}

Almost everything in this manual is a consequence of three decisions, and none
of the rest makes much sense without them.

\subsection{The flux is a solved unknown}

What experiments and downstream codes want out of an equilibrium code is not
$\psi$ but the magnetic field, and the field is built from $\gradbar\psi$. A
method that computes $\psi$ to order $k+1$ and then differentiates it delivers
the field at order $k$ --- one order worse, in precisely the quantity people
use.

{\meq} therefore treats the scaled gradient as an unknown in its own right.
Introduce the \emph{flux}
%
\begin{equation}
	q := \frac{1}{r}\, \gradbar\psi ,
	\label{eq:flux}
\end{equation}
%
and \eqref{eq:gs} becomes the first-order system
%
\begin{equation}
	q - \frac{1}{r}\gradbar\psi = 0,
	\qquad
	-\gradbar\cdot q = \frac{F}{r},
	\qquad
	\psi = 0 \ \text{ on } \Gamma .
	\label{eq:first-order}
\end{equation}
%
Discretising $q$ directly gives the derivative at the \emph{same} order as the
potential. That is the whole reason to prefer a mixed method for this equation,
and it is not a numerical convenience: it is the deliverable.

It goes on paying off in places nobody designed for, and the pattern recurs
throughout this manual. A critical point of $\psi$ --- the magnetic axis, or an
X-point --- is a \emph{root} of $q$ rather than a turning point of something
differentiated, so it is found to the flux's own order and located inside an
element rather than at a mesh vertex. The weight in every flux-surface average
is $1/|\gradbar\psi| = 1/(r|q|)$, available pointwise and again at the flux's own
order. The boundary condition on a curved plasma boundary is carried inward
along a path by integrating $q$ along it. And where the computational domain
stops short of the plasma boundary, the solution is continued across the gap by
a Taylor step in $q$ rather than by extrapolating a polynomial. In every one of
those, a code that differentiated its potential would be a full order worse.

\subsection{Newton, not Picard}

Somebody has to iterate, because \eqref{eq:gs} is semi-linear. The published
hybridizable discontinuous Galerkin solvers for this equation use
Anderson-accelerated Picard iteration, keeping $F$ as opaque problem data so
that the solver relies only on the discretisation of the operator
\cite{SanchezVizuetSolano2019}. They pay for that by iterating on \emph{every}
source, including one linear in $\psi$ that could have been folded into the
bilinear form.

{\meq} makes the opposite trade and uses Newton, which puts
$\partial F/\partial\psi$ into the operator as a mass term on the potential
block. The support for that choice comes from outside the discontinuous
Galerkin literature. CEDRES++ reports that fixed-point iterations `usually
suffer from very slow convergence or even fail to converge, which made
researchers move towards Newton-type methods', and names vertically unstable
plasmas as a case where Picard does not converge at all \cite{Heumann2015}.
Lackner \cite{Lackner1976} says the same from the other end: plain Picard `will
converge to the physically trivial solution $\psi \equiv 0$ if admitted by the
formulation of the problem'. Serino and co-workers \cite{Serino2024} built a
Newton solver for this equation in the same finite element library precisely
because conventional Picard-based solvers failed to converge on the cases they
cared about.

The obligation this creates reaches every corner of the code and is stated here
because it constrains what a user may write:

\begin{quote}
	\textbf{Every source must be able to differentiate itself with respect to
	$\psi$, and every profile must be differentiable twice.} A profile that
	cannot supply its own derivative cannot be used.
\end{quote}

The second derivative is not decorative. A rotating plasma's source is already
a $\psi$-derivative of something built from flux functions, so its Jacobian
spends a second derivative of every input; no reparametrisation avoids it.

The reason the obligation must be taken seriously, rather than merely noted, is
that a wrong Jacobian is nearly invisible. Newton converges to the same
discrete solution whatever Jacobian carried it there, so an error in
$\partial F/\partial\psi$ leaves every error norm and every convergence rate
untouched. What changes is only the path: the observed order of the Newton
iteration drops from two to one, and the iteration count rises. \emph{A
convergence table cannot see a Jacobian error at all.} That is why {\meq}
prints the observed order of every Newton run, why its test suite differences
the assembled Jacobian against the assembled residual, and why
\chref{ch:embedding} returns to the point when a consumer's own Jacobian is at
stake.

\subsection{Hybridization}

The third commitment is the discretisation. {\meq} implements the LDG-H
hybridizable discontinuous Galerkin method of S\'anchez-Vizuet and Solano
\cite{SanchezVizuetSolano2019}, with the adaptivity, error estimator and
curved-boundary treatment of S\'anchez-Vizuet, Solano and Cerfon
\cite{SanchezVizuet2020adaptive}. Those two papers are not background reading:
they are the method, and \texttt{src/meq} is an implementation of them.

Three finite element spaces carry the three unknowns of \eqref{eq:first-order},
all of the \emph{same} polynomial degree $k$: a discontinuous vector space for
the flux, a discontinuous scalar space for the potential, and a scalar space
living on the faces of the mesh for a new unknown, the \emph{trace}
$\hat\psi_h$, which is the potential's value on a face. Equal order is
admissible because hybridization removes the inf-sup compatibility condition a
classical mixed method would impose between flux and potential
\cite{CockburnGopalakrishnanLazarov2009}, and both $\psi$ and $q$ then converge
at $k+1$.

What hybridization buys is that the flux and the potential can be eliminated
\emph{element by element}, each element's contribution expressed in terms of the
trace on its own faces. What is left is one globally coupled system in the trace
alone, whose size is the number of face degrees of freedom and is independent of
how many volume unknowns each element carries. That is why a high-degree run
costs far less than its unknown count suggests, and it is why the choice of
linear solver in \chref{ch:using} is a choice about one sparse matrix rather
than about the whole discretisation.

What it costs is the subject of \chref{ch:using}'s section on nonlinear solvers,
and it is a real cost: when $F$ depends on $\psi$, the element-local elimination
can itself be nonlinear, and then there is a small nonlinear solve inside every
element for every evaluation of the global residual. Neither of the other two
commitments interacts with the third as sharply as this, and it is the single
largest robustness lever in the code.

\section{Scope}

It is worth being plain about what {\meq} does not do, since each omission is a
deliberate boundary rather than an unfinished corner.

\begin{description}

\item[No free boundary.] The plasma boundary is an input. The complementary
	problem --- in which external coils are given and the plasma finds its own
	shape, with the boundary determined as part of the solution --- is harder,
	and {\meq} does not solve it. Serino and co-workers \cite{Serino2024} note
	that the fixed problem is `significantly easier', which is worth keeping in
	mind when comparing against a free-boundary code. The design anticipates the
	extension in one specific place: when profiles are given in normalised flux
	the flux on axis is already an unknown of the nonlinear system, closed by a
	bordered Newton, and the boundary flux would enter as a second border row
	and column of exactly the same shape.

\item[No parallelism.] {\meq} is serial throughout, and there is no MPI in it.
	This is a judgement about problem size rather than an omission: a
	two-dimensional trace system of a few tens of thousands of unknowns
	factorises comfortably on one core with a sparse direct solver, and at that
	size a direct solve beats anything iterative. The judgement would not
	survive three dimensions.

\item[Two dimensions, and axisymmetry.] The reduction to \eqref{eq:gs} assumes
	it.

\item[No open surfaces.] Everything downstream of the solve --- the tracer, the
	flux-surface averages, the disc representation --- assumes closed surfaces
	that return to their starting point. A surface terminating on the domain
	boundary has endpoints and wants different rules, and that is deferred with
	free boundary.

\item[Nothing near the separatrix.] Approaching an X-point, $1/|\gradbar\psi|$
	diverges, the surface develops a corner and every flux-surface average
	diverges logarithmically with it. That is physics rather than a defect, but
	it means a consumer must decide where to cut, and record the decision.

\end{description}

\section{What it is built on}

The finite element machinery is MFEM \cite{Anderson2021}, and specifically its
mixed \texttt{DarcyForm} and the hybridizable discontinuous Galerkin integrators
that go with it. {\meq} is an early user of that work and requires a development
branch of the library rather than a release; \chref{ch:using} says which, and
why the requirement cannot be relaxed.

The curved-boundary treatment is the transfer technique of Cockburn and Solano
\cite{CockburnSolano2012}, in which the discrete problem is posed on a polygonal
subdomain lying strictly inside the true domain and the boundary condition is
transferred outward along short paths. Its essential property --- and the reason
it is more than a heuristic --- is that optimal order survives even when the gap
between the computational boundary and the true one is only $O(h)$. Earlier
methods of this kind needed the gap to shrink like $h^{k+1}$, which in practice
means fitting the mesh after all.

The distinctive combination is \emph{a mixed high-order method, on a mesh that
does not conform to the plasma boundary, driven by Newton}. Each of those three
is well established on its own. Taking all three at once is what produces most
of the interesting behaviour documented here, and most of the traps.

\section{How to read this manual}

\Chref{ch:using} is for somebody who wants to run {\meq}: building it, the one
hard dependency, the configuration file, the three output formats and why there
are three of them, and the options the code exposes together with the standing
advice about each. \Chref{ch:embedding} is for somebody who wants to call
{\meq} from a larger code --- a transport solver, an optimiser, a design loop
--- and it is the chapter with the most in it that appears nowhere else: the
solver object and its lifecycle, sources and profiles as the physics input, warm
starting between nearby equilibria, and the flux-surface machinery that turns a
solved equilibrium into one-dimensional geometry. Appendix~\ref{ap:config} is
the complete syntax of a configuration file.

A note on numbers, which applies throughout. Many of {\meq}'s design choices
could only be settled by measurement, and this manual says so where that is the
case. It does not reproduce the measurements. A number in a manual goes stale
silently; the numbers therefore live in the source tree beside the tests that
produce them. Where a choice is exposed as an option --- the polynomial degree,
the trace solver, the assembly mode, the nonlinear ordering, the globalisation,
the number of threads --- the right setting depends on the problem, the mesh and
the machine, and the only way to find it is to measure it there. Treat every
recommendation here as the default it is.
